\documentclass[10pt,a4paper]{amsart}
\usepackage[margin=19mm]{geometry}
\usepackage{amsmath,amssymb,mathtools}
\usepackage[scr=rsfs,cal=euler]{mathalfa}
\usepackage{xfrac}
\usepackage[unicode,hidelinks,bookmarks=false]{hyperref}
\newcommand{\ie}{i.e.\ }
\newcommand{\cf}{cf.\ }
\newcommand{\resp}{respectively\ }
\newcommand{\Subsection}{\subsection}
\providecommand{\nobreakdash}{\nobreak\hbox{-}}
\setlength{\emergencystretch}{2em}
\multlinegap0pt
\usepackage{bm}
\usepackage{bbm}
\usepackage{dsfont}
\usepackage{xspace}
\usepackage{cancel}
\usepackage{mathtools}
\usepackage{amsthm}
\makeatletter
\@ifundefined{theorem}{\newtheorem{theorem}{Theorem}}{}
\@ifundefined{claim}{\newtheorem{claim}{Claim}}{}
\@ifundefined{lemma}{\newtheorem{lemma}[theorem]{Lemma}}{}
\@ifundefined{proposition}{\newtheorem{proposition}[theorem]{Proposition}}{}
\makeatother
\DeclareMathOperator{\Diag}{Diag}
\DeclareMathOperator{\Sym}{Sym}
\newcommand{\N}{\mathbb{N}}
\title[The distinguishing number of complete bipartite and crown gr]{The distinguishing number of complete bipartite and crown graphs}
\author{Lei Chen, Alice Devillers, Luke Morgan, Friedrich Rober}
\date{Source: arXiv:2601.15913v1 (CC BY 4.0)}
\begin{document}
\maketitle

\noindent\emph{Source and licence.} The distinguishing number of complete bipartite and crown graphs. Version arXiv:2601.15913v1 (CC BY 4.0), distributed under the Creative Commons Attribution 4.0 International licence, as recorded in the arXiv licence metadata for this exact version and shown on the abstract page https://arxiv.org/abs/2601.15913v1. Full rights notice: licenses/arxiv-2601.15913-CC-BY-4.0.txt.\par\smallskip

\noindent\textit{Editorial scope.} Counted unit: the Proposition whose source label is \texttt{prop:DGdiag-Sym} in the article (source0.tex lines 1100--1104, source label \texttt{prop:DGdiag-Sym}), delivered together with the article's own complete original proof of it (source0.tex lines 1105--1276, one \texttt{proof} environment of 172 lines). No mathematics is written, repaired or re-derived: statement and proof are the article's own text line for line. Required context retained uncounted: V is partition (lines 1056--1070); lem:partsaresubsets (lines 1072--1079). Every definition, display and label of those lines is retained unchanged. Omitted from this package and not redistributed: 1--1055, 1071--1071, 1080--1099, 1277--1526. Nothing inside a retained excerpt is omitted or altered: every retained body is delivered exactly as the article's lines are written, so no omission receipt of any kind accompanies this package. Citations: the retained excerpts carry no citation command at all, so no bibliography entry of the article is transcribed and no reference is redistributed. Reference adaptation: none was needed and none was made; every reference inside the delivered unit resolves inside it, so no unresolved reference or citation is produced. Printed slips kept literal: no printed word, symbol, hypothesis or number of the counted statement or of its proof was altered. Layout and class: the article's own class is \texttt{amsart} with packages including booktabs, a4wide, amsmath, amsthm, amssymb, amsfonts, mathtools, geometry, xcolor, graphicx, hyperref, enumitem, calc, tikz; the wrapper is the lane's pinned \texttt{amsart} editorial wrapper at 10pt on A4, which restates the theorem-like environments the retained text needs and adds the article's own macro definitions that the retained text needs (\texttt{\textbackslash Diag}, \texttt{\textbackslash N}, \texttt{\textbackslash Sym}), with extra wrapper packages (bm, bbm, dsfont, xspace, cancel, mathtools). The delivered numbering is the unit's own. Assets: no retained span contains a figure environment, a graphics-inclusion command, a PDF-page inclusion, a table or a TikZ picture; no image file is copied into this package, the package declares no asset dependency and no image file is redistributed with it. Licence: version arXiv:2601.15913v1 of this article is distributed under the Creative Commons Attribution 4.0 International licence, as recorded in the arXiv licence metadata for this exact version and shown on the abstract page https://arxiv.org/abs/2601.15913v1. A full rights notice is delivered at licenses/arxiv-2601.15913-CC-BY-4.0.txt. Characters: every retained excerpt is pure ASCII, so the delivered engine, XeLaTeX with its default Latin Modern text font, reports no missing character.
\begin{lemma}
\label{V is partition}
Let \(n,q \in \N\) with $q\leqslant n$.
\begin{enumerate}
    \item[(a)] For $1\leqslant i\leqslant q$, $\mathcal{U}_n(i,q)$ is non-empty, and \(
    \{\mathcal{U}_n(i, q) : 1 \leqslant i \leqslant q\}
\)
is a partition of $\Delta'$.
\item[(b)]  For any $i$ such that \((i-1)q < n\), $\mathcal{V}_n(i,q)$ is non-empty, and 
\(
    \{\mathcal{V}_n(i, q) : 1 \leqslant i \leqslant d\}
\)
is a partition of $\Delta$ if \((d-1)q<n\leqslant d \cdot q \).
\end{enumerate}
\end{lemma}

\begin{lemma}\label{lem:partsaresubsets}  Let \(n,q \in \N\) with $q\leqslant n$. Let $\Gamma=K_{n,n}-nK_2$ and assume  $G^+= \Diag(H\times H)$ where $H\leq\Sym(n)$. 

Suppose \(\mathcal{P}:=\{P_{1},\ldots, P_{m}\}\) is a partition of the vertex-set of $\Gamma$ %=K_{n,n}-nK_2$ 
such that for each $i$
\[P_i\cap\Delta\text{ is a subset of }\mathcal{V}_n(a, q)\text{ for some }a, \]
\[P_i\cap\Delta'\text{ is a subset of }\mathcal{U}_n(b, q)\text{ for some }b. \]
Then $G^+_{(\mathcal{P})}$ is trivial.    
\end{lemma}

\begin{proposition}\label{prop:DGdiag-Sym}
    Let \(\Gamma:=K_{n,n}-n K_2\), where $n\geqslant 3$, and assume that $G^+=\Diag(\Sym(n)\times \Sym(n))$.
    %\(G:=\Aut(\Gamma)=\langle  \Diag(\Sym(n)\times \Sym(n)),\tau\rangle\).
    Then \(D(G)=\lfloor \sqrt{n}\rfloor+1\).
\end{proposition} 

\begin{proof}
 
Let \(k:=D(G)\) and $\Sigma:=\{\Sigma_1 , \Sigma_2, \ldots,\Sigma_k\}$  be a distinguishing  partition of \(V\Gamma\), that is, \(G_{(\Sigma)}=\bigcap_{i=1}^{k}G_{\Sigma_{i}}=1\).
\begin{claim}\label{u_idiffparts}
For any \(1\leqslant i\leqslant  k\), any two vertices of \(\{u_j \mid v_j\in \Sigma_i\}\) lie in two different parts of \(\Sigma\).
\end{claim}
Indeed, otherwise, there would exist \(\Sigma_{i'}\in\Sigma\) and \(u_s, u_t\in \Sigma_{i'}\) with \(v_s,v_t\in \Sigma_i\), and the map \(\pi\coloneqq ((s,t),(s,t))\in G^+\) stabilises each element of  \(\Sigma\).
%Then the transposition \(\tau:=(s,t)\) stabilises \(\Sigma_{1}\) and \(\Sigma_{i}\). 
%Moreover, \(\tau\in \Sigma_{j}\) for \(j\neq 1,i\) since \(\tau\) stabilises \(\{1,\ldots,n\}\setminus\{1,i\}\).
Thus \(\pi\in G_{(\Sigma)}\), which is a contradiction, proving the claim.

By Claim \ref{u_idiffparts}, it follows that \(|\Sigma_{i}\cap \Delta|\leqslant k\) for all \(i\).
Thus \(n=|\Delta|=\sum_{i=1}^{k}|\Sigma_{i}\cap\Delta|\leqslant k^{2}\).
Note that \(k\) is an integer, which implies that \(k\geqslant \lceil\sqrt{n}\,\rceil.\)
%\(k\geqslant \lfloor \sqrt{n}\rfloor+1\) when \(n\) is not a square.



First we deal with the case where \(n\) is not a square. 
Let \(\ell:=\lceil\sqrt{n}\,\rceil=\lfloor \sqrt{n}\rfloor+1\).
Then \(\ell-1\) is the largest integer such that its square is strictly less than \(n\),
thus \((\ell-1)^{2} < n < \ell^2\).
In particular, \(n=(\ell-1)^{2}+p\), where \(1 \leqslant p \leqslant 2\ell-2\).  We show that \(k\leqslant \ell\) by displaying a partition \(\mathcal{P}\) of size \(\ell\) with trivial stabiliser \(G_{(\mathcal{P})}\). This in turn will show that \(k=\ell\) when \(n\) is not a square.
We split the analysis into two cases: (i) \(1 \leqslant p< \ell\); (ii) \(\ell \leqslant p \leqslant 2\ell-2\).


 \medskip\noindent
\emph{Case (i):  \(1 \leqslant p< \ell\)}. \\Let \(\mathcal{P}:=\{P_{1},\ldots, P_{\ell}\}\) where
\begin{equation*} 
P_{i}  \coloneqq
\begin{cases}
   \mathcal{V}_n(i,\ell-1) \cup \mathcal{U}_n(i,\ell-1),& \text{if } 1\leqslant i < \ell;\\ 
   \mathcal{V}_n(i,\ell-1), & \text{if } i=\ell.
\end{cases}
\end{equation*}

This is a partition of size $\ell$ of \(V\Gamma\),
since
\(
\{P_i \cap \Delta' : 1 \leq i \leq \ell\}
= \{\mathcal{U}_n(i, \ell-1) : 1 \leq i \leq \ell - 1\} \cup \{\varnothing\}
\)
is a partition of $\Delta'$ by Lemma \ref{V is partition}(a)
and \(
\{P_i \cap \Delta : 1 \leq i \leq \ell\}
= \{\mathcal{V}_n(i, \ell-1) : 1 \leq i \leq \ell\}
\)
is a partition of $\Delta$
by applying Lemma \ref{V is partition}(b) with $d = \ell$ and $q = \ell - 1$, as $(\ell-1)^2 < n = (\ell-1)^{2} + p \leqslant  \ell \, (\ell - 1).$

% \vspace*{2em}
%\subfile{../tikz/crown_graph/Sn/case1}

%\begin{claim}\label{partstab}
 %Let \(g\in G_{(\mathcal{P})}\). Then $g$ fixes $v_s$ for all $s\leqslant n$.
 %\end{claim}


Note that \(P_{\ell}\subseteq\Delta\) and so $G_{(\mathcal{P})}=G^+_{(\mathcal{P})}$. 
Observe that \(\mathcal{P}\) satisfies the conditions of Lemma \ref{lem:partsaresubsets}. It follows that $G^+_{(\mathcal{P})}$ is trivial.
Thus $\mathcal{P}$ is a distinguishing partition for $G$.
%This implies that \(g\in  G^{+}= \Diag(\Sym(n)\times \Sym(n))\) and so $g=(h,h)$ for some $h\in \Sym(n)$. 
% 
% Suppose $h$ does not fix $s$.  Since \(g\in G_{(\mathcal{P})}\), $v_s^g=v_{s^h}\in P_i$. In particular $0<|s-s^h|<\ell-1$. Then $u_s^g=u_{s^h}$ and $u_s$ are also in the same part $P_j$ and therefore $s\equiv s^h\equiv j \;(\bmod\; \ell - 1)$, contradicting $0<|s-s^h|<\ell-1$. It follows that $h=\mathrm{id}$, and so $g=(\mathrm{id},\mathrm{id}).$
 
%Therefore \(G_{(\mathcal{P})}\) is trivial.

\medskip\noindent
\emph{Case (ii): \(\ell\leqslant p\leqslant 2\ell-2\)}.\\ Let \(\mathcal{P}:=\{P_{1},\ldots, P_{\ell}\}\) where
\begin{equation*}
    P_{i}:=\begin{cases}
       \mathcal{V}_n(i,\ell)\cup \mathcal{U}_n(i,\ell)&\text{if }1\leqslant i\leqslant \ell-2\\
        
        \mathcal{V}_n(i,\ell)\cup  \mathcal{U}_n(\ell,\ell)&\text{if } i = \ell-1\\
        
        \mathcal{V}_n(i,\ell)\cup  \mathcal{U}_n(\ell-1,\ell)&\text{if }i=\ell.
    \end{cases}
\end{equation*}

Arguing similarly to the previous case, we see that this is a partition of  \(V\Gamma\) of size $\ell$,
by applying Lemma \ref{V is partition}(b) with $d = \ell$ and $q = \ell$, as $(\ell-1)\ell<n=(\ell-1)^{2}+p\leqslant \ell^2.$

Note that \(|P_{\ell - 1}\cap\Delta'|=\ell-1<\ell=|P_{\ell - 1}\cap\Delta|\),
and so $G_{(\mathcal{P})}=G^+_{(\mathcal{P})}$. 
Observe that \(\mathcal{P}\) satisfies the conditions of Lemma \ref{lem:partsaresubsets}. It follows that $G^+_{(\mathcal{P})}$ is trivial.
Thus $\mathcal{P}$ is a distinguishing partition for $G$.

% \vspace*{2em}
%\subfile{../tikz/crown_graph/Sn/case2}
%We now show that $G_{(\mathcal{P})}$ is trivial. 
 %Let \(g\in G_{(\mathcal{P})}\).
%Observe that \(|P_{\ell - 1}\cap\Delta'|=\ell-1<\ell=|P_{\ell - 1}\cap\Delta|\). Since $g$ preserves $P_{\ell - 1}$, it follows that \(g=(h,h)\in G^{+}\).
%
%Suppose $h$ does not fix $s$.  Since \(g\in G_{(\mathcal{P})}\), we get that $0<|s-s^h|<\ell$. Then $u_s^g=u_{s^h}$ and $u_s$ are also in the same part $P_j$ and therefore $s\equiv s^h \;(\bmod\; \ell )$, contradicting $0<|s-s^h|<\ell$. It follows that $h=\mathrm{id}$, and so $g=(\mathrm{id},\mathrm{id}).$ 
%Therefore \(G_{(\mathcal{P})}\) is trivial.



%With a similar argument used in Case (i) we conclude that \(G_{(\mathcal{P})}\) is trivial.


From these partitions of size \(\ell\)  we see that \(D(G)=\lfloor \sqrt{n}\rfloor+1\) when \(n\) is not a square.
Thus it remains to show that the result also holds for \(n=\ell^{2}\) a square.
Note that \(k^{2}\geqslant n\), and so \(k\geqslant \ell\). 

\medskip
First assume that \(k=\ell\).
Recall that \(|\Sigma_{i}\cap\Delta|\leqslant k=\ell\)
and \(\sum_{i=1}^{\ell}|\Sigma_{i}\cap\Delta|=n=\ell^{2}\).
Hence \(|\Sigma_{i}\cap \Delta|=\ell\) for \(1\leqslant i\leq\ell\). 
Similarly one can show that \(|\Sigma_{i}\cap\Delta'|=\ell\) for \(1\leqslant i\leqslant \ell\).
By Claim \ref{u_idiffparts}, for any \(i\), any two vertices of \(\{u_j|v_j\in \Sigma_i\}\) lie in two different parts of \(\Sigma\), so \(\{u_j \mid v_j\in \Sigma_i\}\) has exactly one vertex from each element of \(\Sigma\).
%
%
%Thus, up to reordering of the vertices, we may assume that \(\Sigma_{i}=\mathcal{V}_n(i,\ell)\cup \mathcal{U}_n(i,\ell)\) for \(1\leqslant i\leqslant \ell\) where \(|\mathcal{V}_n(i,\ell)|=|\mathcal{U}_n(i,\ell)|=\ell\) for each \(i\).
%{\color{purple} Not sure if we want to explain the argument here more rigorously or not. I had to think a few minutes why we could reorder the vertices like that by using a permutation $(x,x) \in \Diag(\Sym(n) \times \Sym(n))$. }
%{\color{red} I am trying something else. Is this clearer or more confusing than reordering?}
Therefore the map $\alpha:\{1,2,\ldots,\ell^2\}\mapsto \{1,2,\ldots,\ell\}^2$ defined by $\alpha(i)=(x,y)$ if $v_i\in \Sigma_x$ and $u_i\in \Sigma_y$, is a bijection.
Let $h$ be the element of $\Sym(n)$, defined by $\alpha(i^h)=(y,x)$ where $\alpha(i)=(x,y)$. Clearly $h$ is an involution.
 Let \(g:=(h,h)\tau\in G\), so that  $v_i^g=(v_{i^h})^\tau=u_{i^h}$ and $u_i^g=(u_{i^h})^\tau=v_{i^h}$. It follows that $g$ is an involution.
 %%%%%{\color{red} here again we need that $\tau \in G$, after I changed the hypothesis slightly. Really need in the preliminary something about $G\setminus G^+$ when $G^+$ is diagonal, see text in blue.}
 
 
 We claim that $g\in G_{({\Sigma})}.$
Let $v_i\in\Sigma_x\cap \Delta$, thus $\alpha(i)=(x,y)$ for some $y$, and so $\alpha(i^h)=(y,x)$, and so $v_i^g=u_{i^h}\in \Sigma_x\cap\Delta'$.
Therefore \((\Sigma_{x}\cap\Delta)^{g}=\Sigma_{x}\cap\Delta'\).
Since \(g \) is an involution we also have that \((\Sigma_{x}\cap\Delta')^{g}=\Sigma_{x}\cap\Delta\).
This shows $\Sigma$ is not a distinguishing partition, a contradiction.  Thus
\(k\geqslant \ell+1\).





%Let \(h\in H\cong \Sym(n)\) mapping \(v_{(i-1)\ell+j}\) to \(v_{(j-1)\ell+i}\) for each \(1\leqslant i,j\leqslant \ell\).  Let \(g:=(h,h)\tau\in G\). Note that \(|h|=2\) and so \(|g|=2\).

%Then  \begin{align*}           
%    (\Sigma_{i}\cap\Delta)^{g}&=\mathcal{V}_n(i,\ell)^g=\{v_{(i-1)\ell+j} \mid  j=1,\ldots, \ell\}^g\\
%&=  \{u_{(i-1)\ell+j} | j=1,\ldots, \ell\}^h  \\
%&=  \{u_{(j-1)\ell+i} \mid j=1,\ldots, \ell\}\\
%&=\mathcal{U}_n(i,\ell)=\Sigma_{i}\cap\Delta'
%\end{align*}
%Since \(g \) is an involution we also have that \((\Sigma_{i}\cap\Delta')^{g}=\Sigma_{i}\cap\Delta\).

%Thus \(g\) stabilises \(\Sigma_{1},\ldots, \Sigma_{\ell}\) and \(g\in G_{(\mathcal{P})}\). 
%This contradicts the  assumption   that $D(G)=k$. Hence \(k\geq\ell+1\).

We will now construct a distinguishing partition with \(\ell+1\) parts, which shows that \(D(G)=\ell+1=\sqrt{n}+1=\lfloor \sqrt{n}\rfloor +1\) when \(n\) is a square.
Consider \(\mathcal{P}:=\{P_{1},\ldots, P_{\ell+1}\}\)  whereby 
\begin{equation}
P_{i}:=
\begin{cases}
    \mathcal{V}_n(i,\ell)\cup\mathcal{U}_n(i,\ell) \setminus \{u_{\ell^{2}}\} & \text{if } 1\leqslant i\leqslant \ell\\
    \{u_{\ell^{2}}\} & \text{if } i=\ell+1.
\end{cases}
\end{equation}
This is clearly a partition of size $\ell+1$ of \(V\Gamma\),
by applying Lemma \ref{V is partition}.


%\vspace*{2em}
%\subfile{../tikz/crown_graph/Sn/case3}


Note that  \(P_{\ell+1}\subseteq \Delta'\),
and so $G_{(\mathcal{P})}=G^+_{(\mathcal{P})}$. 
Observe that \(\mathcal{P}\) satisfies the conditions of Lemma \ref{lem:partsaresubsets}. It follows that $G^+_{(\mathcal{P})}$ is trivial.
Thus $\mathcal{P}$ is a distinguishing partition for $G$.
%Then since \(P_{\ell+1}\subseteq \Delta'\), it follows that \(G_{(\mathcal{P})}\leqslant G_{\Delta'}= G^{+}\).
%Applying a similar argument as in Case (i) we conclude that \(G_{(\mathcal{P})}\)  is trivial. 
%Hence \(D(G)=\ell+1=\lfloor \sqrt{n}\rfloor+1\), proving the result for the case \(n\) square.
\end{proof}

\end{document}
